\documentclass[border=10pt,tikz]{standalone}
\usepackage{fontspec}
\setmainfont{Apple SD Gothic Neo}
\setsansfont{Apple SD Gothic Neo}
\usepackage{tikz}
\usetikzlibrary{positioning,arrows.meta,shapes,fit,backgrounds,calc}
\begin{document}
% _design.tex — shared TikZ design system for all blog diagrams
% Inputs nothing on its own — meant to be \input{} from each diagram .tex
% AFTER \begin{document}, BEFORE \begin{tikzpicture}.
%
% Provides a unified visual language across GoF / DSA / EMC++ / Beautiful series:
%   - Colors: muted, accessible, color-coded by role
%   - Node styles: consistent sizing, rounded, subtle borders
%   - Edge styles: UML-style inheritance, aggregation, plus dashed alternates
%
% Usage:
%   \documentclass[border=10pt,tikz]{standalone}
%   \usepackage{tikz}
%   \usetikzlibrary{positioning,arrows.meta,shapes,fit,backgrounds}
%   \begin{document}
%   \input{../_design.tex}     % adjust relative path
%   \begin{tikzpicture}[blog]  % apply 'blog' style
%   ...

\definecolor{absfill}{RGB}{220,232,247}   % cool blue
\definecolor{absbord}{RGB}{72,118,182}
\definecolor{confill}{RGB}{249,222,222}   % warm red
\definecolor{conbord}{RGB}{197,93,93}
\definecolor{impfill}{RGB}{223,238,222}   % soft green
\definecolor{impbord}{RGB}{99,154,93}
\definecolor{clifill}{RGB}{251,239,205}   % pale yellow
\definecolor{clibord}{RGB}{200,167,72}
\definecolor{hifill} {RGB}{253,224,200}   % accent orange
\definecolor{hibord} {RGB}{210,135,68}
\definecolor{nodefill}{RGB}{250,250,250}
\definecolor{nodebord}{RGB}{120,120,120}
\definecolor{edgec}  {RGB}{80,80,80}
\definecolor{edgesc} {RGB}{145,145,145}
\definecolor{groupbord}{RGB}{200,200,200}

\tikzset{
  blog/.style={
    font=\sffamily\small,
    every node/.style={font=\sffamily\small}
  },
  % --- node roles ---
  cls/.style={
    rectangle, rounded corners=2pt,
    draw=nodebord, line width=0.4pt,
    minimum height=8mm, minimum width=28mm,
    inner sep=3pt, fill=nodefill, align=center
  },
  abs/.style={cls, draw=absbord, fill=absfill, font=\sffamily\itshape\small},
  con/.style={cls, draw=conbord, fill=confill},
  imp/.style={cls, draw=impbord, fill=impfill},
  cli/.style={cls, draw=clibord, fill=clifill},
  hi/.style={cls, draw=hibord, fill=hifill, font=\sffamily\bfseries\small},
  % --- circle (for trees / graph nodes) ---
  vx/.style={
    circle, draw=absbord, line width=0.4pt,
    minimum size=8mm, fill=absfill, inner sep=0pt
  },
  vxc/.style={vx, draw=conbord, fill=confill},
  vxh/.style={vx, draw=hibord, fill=hifill, font=\sffamily\bfseries},
  % --- decision diamond ---
  q/.style={
    diamond, draw=clibord, line width=0.4pt, fill=clifill,
    inner sep=2pt, aspect=2.4, align=center
  },
  % --- edges ---
  inh/.style={                                  % UML inheritance
    -{Triangle[length=3mm,width=3mm,open]},
    draw=edgec, line width=0.5pt
  },
  agg/.style={                                  % UML aggregation
    {Diamond[length=2.5mm,width=2.5mm,open]}-{Stealth[length=2.2mm]},
    draw=edgec, line width=0.45pt
  },
  uses/.style={                                 % directed reference
    -{Stealth[length=2.2mm]}, draw=edgec, line width=0.5pt
  },
  arr/.style={                                  % plain arrow (graphs/trees)
    -{Stealth[length=2mm]}, draw=edgec, line width=0.45pt
  },
  edge/.style={                                 % undirected (trees)
    draw=edgec, line width=0.45pt
  },
  alt/.style={                                  % dashed: similar/alt
    -{Stealth[length=2mm]}, dashed, draw=edgesc, line width=0.45pt
  },
  ret/.style={                                  % dashed creates/returns
    -{Stealth[length=2mm]}, dashed, draw=edgesc, line width=0.45pt
  },
  thr/.style={                                  % threading link (red dashed)
    ->, dashed, draw=conbord, line width=0.5pt
  },
  % --- group / region ---
  group/.style={
    rectangle, rounded corners=4pt,
    draw=groupbord, line width=0.4pt,
    inner sep=10pt, fill=black!2
  },
  glabel/.style={font=\sffamily\bfseries\small, anchor=north west, inner sep=2pt}
}

% Korean font convenience macro — included by every Hangul diagram
\providecommand{\KOR}{}

\begin{tikzpicture}[blog,
  corebox/.style={rectangle, rounded corners=3pt, draw=absbord, line width=0.5pt, minimum width=46mm, minimum height=12mm, inner sep=4pt, fill=absfill, font=\sffamily\bfseries\footnotesize, align=center},
  cachebox/.style={rectangle, draw=clibord, line width=0.5pt, minimum width=58mm, minimum height=22mm, inner sep=4pt, fill=clifill, align=center},
  membox/.style={rectangle, draw=hibord, line width=0.5pt, minimum width=86mm, minimum height=14mm, inner sep=4pt, fill=hifill, align=center},
  varA/.style={rectangle, draw=impbord, line width=0.4pt, fill=impfill, minimum width=14mm, minimum height=7mm, inner sep=1pt, font=\sffamily\bfseries\footnotesize},
  varB/.style={rectangle, draw=conbord, line width=0.4pt, fill=confill, minimum width=14mm, minimum height=7mm, inner sep=1pt, font=\sffamily\bfseries\footnotesize},
  blank/.style={rectangle, draw=nodebord, line width=0.3pt, fill=nodefill, minimum width=14mm, minimum height=7mm, inner sep=1pt, font=\sffamily\footnotesize}
]

% Title
\node[font=\sffamily\bfseries\small, text=inkc] at (0, 5.6) {False Sharing — 64-byte line을 가정한 예};

% Core 0 (left) and Core 1 (right)
\node[corebox] (c0) at (-6, 4.2) {Core 0\\[2pt] writes \texttt{a}};
\node[corebox] (c1) at ( 6, 4.2) {Core 1\\[2pt] writes \texttt{b}};

% Cache on each core — label outside the box
\node[font=\sffamily\bfseries\footnotesize, text=inkc, anchor=west] at (-8.9, 3.25) {L1 Cache (Core 0)};
\node[cachebox] (cache0) at (-6, 1.85) {};
\node[varA] at (-7.5, 1.85) {a};
\node[varB] at (-6.0, 1.85) {b};
\node[blank] at (-4.5, 1.85) {...};

\node[font=\sffamily\bfseries\footnotesize, text=inkc, anchor=west] at (3.1, 3.25) {L1 Cache (Core 1)};
\node[cachebox] (cache1) at ( 6, 1.85) {};
\node[varA] at (4.5, 1.85) {a};
\node[varB] at (6.0, 1.85) {b};
\node[blank] at (7.5, 1.85) {...};

\draw[arr] (c0.south) -- (cache0.north);
\draw[arr] (c1.south) -- (cache1.north);

% State labels
\node[font=\sffamily\footnotesize, text=conbord] at (-6, 0.0) {MESI: M $\leftrightarrow$ I 반복};
\node[font=\sffamily\footnotesize, text=conbord] at ( 6, 0.0) {MESI: I $\leftrightarrow$ M 반복};

% Coherence traffic — ping pong
\draw[focus, draw=warnbord, line width=0.9pt, <->] (cache0.east) .. controls (-2, 2.4) and (2, 2.4) .. (cache1.west);
\node[font=\sffamily\bfseries\footnotesize, text=warnbord, align=center, fill=white, inner sep=2pt] at (0, 2.6) {coherence \\[2pt] ping-pong};

% Shared memory — single 64-byte line
\node[font=\sffamily\bfseries\footnotesize, text=inkc] at (0, -0.6) {64-byte cache line in DRAM};
\node[membox] (mem) at (0, -1.6) {};
\node[varA] at (-3.5, -1.6) {a};
\node[varB] at (-2.0, -1.6) {b};
\node[blank, minimum width=46mm] at (1.8, -1.6) {... padding / other vars ...};

\draw[arr, draw=mutec] (cache0.south) -- ($(cache0.south)+(0,-0.4)$) -| (mem.north -| -3.5, 0);
\draw[arr, draw=mutec] (cache1.south) -- ($(cache1.south)+(0,-0.4)$) -| (mem.north -| 2.0, 0);

% Bottom annotation
\node[note, anchor=north, text width=160mm, align=left] at (0, -3.0) {
  \textbf{문제}: \texttt{a}와 \texttt{b}는 서로 다른 변수지만 같은 64-byte line에 속함. \\[2pt]
  서로 다른 field의 동시 접근 중 write가 있으면 ownership 이동과 invalidation이 발생할 수 있음. \\[2pt]
  \textbf{확인}: 실제 cache protocol, field offset, PMU 표본과 workload 결과를 함께 대조.
};
\node[note, anchor=north, fill=okfill, draw=okbord, text width=160mm, align=left] at (0, -5.6) {
  \textbf{분리 예}: alignment와 원소 stride를 함께 확인. 이 가정의 배치 검증은 성능 실측과 별도.
};

\end{tikzpicture}
\end{document}
